\section{项目背景与建设目标}
\label{sec:background}

\subsection{企业多模型接入与混合调用场景}
随着大语言模型（Large Language Models, LLMs）在企业知识管理、代码研发、业务分析与客服自动化等核心场景的深度渗透，现代企业的基础设施架构普遍演进为“多模型混合接入”格局。为了在推理延迟、任务能力、采购成本以及可用性冗余之间取得最优权衡，企业统一 AI 网关通常同时路由并对接多家第三方前沿模型服务商，典型代表包括专注于通用逻辑与复杂工程代码生成的 OpenAI 系列模型、擅长超长文本解析与文档精炼的 Kimi 系列模型，以及具备深度推理与专业代码能力的 DeepSeek 系列模型。

在企业实际生产环境中，员工通过统一身份认证（SSO）访问内部门户或 IDE 插件发起模型调用，系统网关实时沉淀了包含用户标识（\texttt{user\_id}）、请求时间戳（\texttt{timestamp}）、目标供应商（\texttt{provider}）、业务任务类型（\texttt{task}）以及端到端执行状态（\texttt{success}）的高频遥测日志数据（Telemetry Data）。

\subsection{企业内部运营与效能度量诉求}
为了支持企业管理层、研发效能团队、财务采购部门以及数据治理团队的精细化决策，企业迫切需要基于沉淀的遥测日志回答以下关键业务问题：
\begin{enumerate}[label=(\arabic*), leftmargin=2em]
  \item \textbf{技术采纳广度（User Prevalence）：} 全公司整体及各业务线员工对大模型能力的渗透率如何？有多少比例的员工真正将大模型引入了日常工作流？
  \item \textbf{业务结构占比（User-Normalized Share）：} 员工主要利用大模型解决哪一类工作？任务分布是以技术研发为主，还是以文档撰写与数据分析为主？
  \item \textbf{交互调用规模（Interaction Volume）：} 整体交互规模与各供应商的承载负荷是多少？各部门的算力需求呈现何种演进趋势？
  \item \textbf{服务完成质量（Success Rate）：} 不同模型供应商在特定任务类型下的端到端成功率如何？是否存在某供应商在特定业务域质量显著劣化的情况？
  \item \textbf{核心热点分布（Top-K Tasks）：} 在长周期演进中，哪些任务分类处于调用前列，值得投入自研工具链进行深度针对性优化？
\end{enumerate}

上述度量结果直接关乎算力预算分配、模型采购合同谈判权重倾斜，以及各部门 AI 辅助效能的客观评估。

\subsection{对外数据共享面临的敏感信息暴露风险}
在企业与大模型服务商的商业合作中，供应商通常要求企业反馈阶段性调用质量与业务渗透统计，以便协同调优模型、优化推理专线并制定后续支持策略。同时，行业协会、外部审计机构或科研合作方也存在对脱敏聚合统计的合法共享诉求。

然而，企业在商业协作中不能直接向外部不可信实体提供明细日志，甚至\textbf{直接发布未经隐私扰动的精确统计值同样蕴含严重的数据安全红线}：
\begin{itemize}[leftmargin=2em]
  \item \textbf{稀疏单元重构攻击（Reconstruction via Sparsity）：} 在按部门、供应商与任务交叉细分的统计表格中，部分特定业务单元（例如“某边缘任务 $\times$ 某模型”）的调用量可能仅由极少数关键员工产生。外部接收方通过分析精确数值，极易重构特定核心研发人员或小微项目组的工作轨迹。
  \item \textbf{差分推断攻击（Differencing Attack）：} 若系统持续对外发布高精度的每日或每周统计值，外部攻击者通过对相邻时间窗口的统计差值进行代数求解，能够高概率推断出某特定员工是否休假、离职或在特定敏感时段发起了密集异常调用。
  \item \textbf{商业竞争与业务架构刺探：} 精确的业务任务比例直接暴露了企业的技术研发重心、正在推进的核心业务形态，乃至企业对不同供应商调用配比的底层商业机密。
\end{itemize}

传统的去标识化方案（如移除员工姓名或对工号做单向 Hash 处理）仅阻断了浅层的身份关联，在面对拥有外部背景知识（Auxiliary Knowledge）的强攻击者时，无法抵御基于行为时序与调用特征的重识别（Re-identification）。

\subsection{系统建设目标}
针对上述业务冲突与安全合规痛点，本项目研发了 \textbf{Private LLM Analytics}（企业大模型使用数据差分隐私分析系统）。系统的核心目标概括如下：
\begin{techbox}{系统总体目标}
构建一套自主受控的企业级差分隐私统计分析基础设施，在设定的严格用户级邻接关系与 $(\varepsilon, \delta)$ 隐私参数下，为企业大模型使用日志生成具备严格数学稳定性保证的聚合统计报告；使企业既能在内部获得高保真度的业务效能度量依据，又能在无需担忧单名员工行为轨迹被外部逆向反推的前提下，安全向外部大模型供应商反馈高质量统计洞察。
\end{techbox}

\subsection{系统核心设计原则}
为了达成上述目标，系统在底层架构设计上严格遵循以下四项原则：
\begin{enumerate}[label=\textbf{原则\arabic*：}, leftmargin=4.5em]
  \item \textbf{指标专属机制（Statistic-Specific）：} 摒弃将所有业务指标简单套用单一通用敏感度与拉普拉斯/高斯加噪的做法。针对比例、占比、总量、成功率与排序等不同统计算子的取值空间，分别推导针对性的贡献限制算子、紧致敏感度与专有联合发布机制。
  \item \textbf{两层误差解耦（Two-Tier Error Decomposition）：} 在效用分析中，清晰界定“贡献截断偏差（Clipping Bias）”与“差分隐私高斯加噪误差（DP Noise Error）”，为工程参数调优与算法演进提供客观的数值归因。
  \item \textbf{规范重放防护（Canonical Replay Protection）：} 基于全维度语义哈希实现严格的幂等发布控制，相同的查询请求强制命中历史发布结果，净增隐私消耗恒等于零，避免生成独立重复噪声样本，防范攻击者通过重复查询求平均来抵消随机噪声的均值恢复攻击。
  \item \textbf{故障闭环安全（Fail-Closed Policy）：} 任何超出预设全生命周期预算上限（Lifetime Budget Cap）、违反接收方作用域策略或存在参数异常的查询，系统均无条件阻断执行，绝不向不可信接收方降级回退原始数据。
\end{enumerate}
